\documentclass[aspectratio=169, lualatex, handout]{beamer}
\makeatletter\def\input@path{{theme/}}\makeatother\usetheme{cipher}

\title{Applied Cryptography}
\author{Nadim Kobeissi}
\institute{American University of Beirut}
\instituteimage{images/aub_white.png}
\date{\today}
\coversubtitle{CMPS 297AD/396AI}
\coverpartname{Part 1: Provable Security}
\coversessionname{1.6: Collision-Resistant Hash Functions}
\coverwebsite{https://appliedcryptography.page}

\begin{document}
\begin{frame}[plain]
	\titlepage
\end{frame}

\begin{frame}{Symmetric primitive example: hash functions}{Reminder}
	\begin{columns}[c]
		\begin{column}{0.55\textwidth}
			\definitionbox{Hash Function Properties}{
				\begin{itemize}\item Takes input of \textbf{any size}[<+->]
					\item Produces output of \textbf{fixed size}
					\item Is \textbf{deterministic} (same input $\rightarrow$ same output)
					\item Even a \textbf{tiny change} in input creates completely different output
					\item Is \textbf{efficient} to compute\end{itemize}
			}
		\end{column}
		\begin{column}{0.45\textwidth}
			\begin{tcolorbox}
				[colback=black!5!white,colframe=ciphergray] $\mathsf{SHA256}(\texttt{hello}) =$ \\ \texttt{2cf24dba5fb0a30e26e83b2ac5}\\ \texttt{b9e29e1b161e5c1fa7425e7304}\\
				\texttt{3362938b9824}

				$\mathsf{SHA256}(\texttt{hullo}) =$ \\ \texttt{7835066a1457504217688c8f5d}\\
				\texttt{06909c6591e0ca78c254ccf174}\\ \texttt{50d0d999cab0}
			\end{tcolorbox}
			\textcolor{cipherprimary}{\textbf{Note:} \small One character change $\rightarrow$
				completely different hash!}
		\end{column}
	\end{columns}
\end{frame}

\begin{frame}{Expected properties of a hash function}{Reminder}
	\begin{columns}[c]
		\begin{column}{0.6\textwidth}
			\begin{itemize}[<+->]
				\item \textbf{Collision resistance}: computationally infeasible to find
				      two different inputs producing the same hash.
				\item \textbf{Preimage resistance}: given the output of a hash function,
				      it is computationally infeasible to reconstruct the original input.
				\item \textbf{Second preimage resistance}: given an input and an output,
				      it's computationally infeasible to find another different input
				      producing the same output.
			\end{itemize}
		\end{column}
		\begin{column}{0.4\textwidth}
			\imagewithcaption{sha2.png}{SHA-2 compression function. Source: Wikipedia}
		\end{column}
	\end{columns}
\end{frame}

\begin{frame}{Hash functions: what are they good for?}{Reminder}
	\begin{itemize}[<+->]
		\item \textbf{Password storage}: Store the hash of the password on the server,
		      not the password itself. Then check candidate passwords against the hash.
		\item \textbf{Data integrity verification}: Hash a file. Later hash it
		      again and compare hashes to check if the file has changed, suffered storage
		      degradation, etc.
		\item \textbf{Proof of work}: Server asks client to hash something a lot of
		      times before they can access some resource. Useful for anti-spam, Bitcoin
		      mining, etc.
	\end{itemize}
\end{frame}

\begin{frame}{Slides not complete}
	\begin{itemize}
		\item This slide deck is not finished and is missing important material. Do not rely on it yet.
	\end{itemize}
\end{frame}

\begin{frame}[plain]
	\titlepage
\end{frame}
\end{document}
